\subsection*{Centering and spectral decompositions}
\begin{problem}[Centering and eigenvalue dispersion]\label{ex:12:centering}
Let $C=I_n-\mathbf1\mathbf1^\top/n$. Prove that $C$ is symmetric
idempotent, that $(Cx)_j=x_j-\bar x$, and that $C\mathbf1=0$.
For symmetric positive semidefinite $A$, prove
$\tr(A^2)-(\tr A)^2/n\geqslant0$.
\end{problem}
\begin{solution}
The matrix $J=\mathbf1\mathbf1^\top/n$ is symmetric and
$J^2=\mathbf1(\mathbf1^\top\mathbf1)\mathbf1^\top/n^2=J$.
Thus $C^\top=C$ and $C^2=I-2J+J^2=C$.
Since $Jx=\bar x\mathbf1$, the coordinate claim follows, and putting
$x=\mathbf1$ gives $C\mathbf1=0$.
If $\lambda=(\lambda_1,\ldots,\lambda_n)^\top$ lists the eigenvalues
of $A$, spectral diagonalization gives
\[
\tr(A^2)-\frac{(\tr A)^2}{n}
=\sum_i(\lambda_i-\bar\lambda)^2
=\lambda^\top C\lambda\geqslant0.
\]
Equality holds exactly when all eigenvalues coincide, equivalently
$A=\bar\lambda I$. The same inequality in fact holds for every real
symmetric $A$.
\end{solution}

\begin{problem}[A resolution of the identity]
Let $P_1,\ldots,P_m$ be symmetric idempotent $n\times n$ matrices with
$\sum_iP_i=I$ and ranks $r_i$. Prove $\sum_i r_i=n$,
$P_iP_j=0$ for $i\ne j$, and that $P_i+P_j$ is symmetric idempotent.
\end{problem}
\begin{solution}
Since the trace of each projector equals its rank,
$\sum_i r_i=\tr(\sum_iP_i)=n$.
For $x\in\operatorname{col}(P_j)$, $P_jx=x$ and
\[
\norm x^2=x^\top\left(\sum_iP_i\right)x
=\norm x^2+\sum_{i\ne j}\norm{P_ix}^2.
\]
Every summand in the remaining sum is nonnegative, so $P_ix=0$ for
$i\ne j$. Apply this to $x=P_jy$ for arbitrary $y$ to get $P_iP_j=0$.
Transposition also gives $P_jP_i=0$.
Thus $(P_i+P_j)^2=P_i+P_j$, and symmetry is inherited from the two
summands. This proves all assertions without a separate rank-inequality
criterion for idempotence.
\end{solution}

\begin{problem}[Weighted sums of perpendicular projectors]
Under the preceding hypotheses, let $A=\sum_i a_iP_i$.
Find its eigenvalues with multiplicities, its determinant, and its inverse
when every $a_i\ne0$.
\end{problem}
\begin{solution}
On $\operatorname{col}(P_i)$ the matrix acts as $a_iI$, since
$AP_i=a_iP_i$. The ranges are perpendicular and their direct sum is
$\R^n$, as $x=\sum_iP_ix$ for every $x$.
Choose an orthonormal basis in each range and concatenate them. The
result diagonalizes $A$ with $a_i$ repeated $r_i$ times; if several
$a_i$ coincide their multiplicities add. Ranges of rank zero contribute
no eigenvectors. Thus $\det A=\prod_i a_i^{r_i}$, with zero-rank factors
omitted. If all $a_i\ne0$, then
\[
A\left(\sum_i a_i^{-1}P_i\right)=\sum_iP_i=I,
\]
and the reversed product is also $I$, so
$A^{-1}=\sum_i a_i^{-1}P_i$.
\end{solution}

\begin{problem}[Equicorrelation]\label{ex:12:equicorrelation}
For $n\geqslant2$, let $A$ have diagonal entries $1$ and off-diagonal
entries $t$. Find its eigenvalues/eigenvectors, determinant, inverse
when it exists, and the exact range of $t$ for positive definiteness.
\end{problem}
\begin{solution}
With $J=\mathbf1\mathbf1^\top/n$ and $C=I-J$,
\[
A=(1-t)I+t\mathbf1\mathbf1^\top
=[1+(n-1)t]J+(1-t)C.
\]
Its eigenvalue on $\R\mathbf1$ is $a=1+(n-1)t$; on
$\mathbf1^\perp$ it is $b=1-t$, of multiplicity $n-1$.
For $t\ne0$ these two eigenvalues are distinct, so their eigenspaces
are exactly those two subspaces. For $t=0$, $A=I$ and every nonzero
vector is an eigenvector. An orthonormal eigenbasis is always
$\mathbf1/\sqrt n$ followed by any orthonormal basis of $\mathbf1^\perp$.
The determinant is $ab^{n-1}$. It is invertible exactly when
$t\notin\{1,-1/(n-1)\}$, with $A^{-1}=a^{-1}J+b^{-1}C$.
Positive definiteness requires both $a,b>0$, hence
$-1/(n-1)<t<1$. At the two endpoints it is positive semidefinite and
singular, with the zero eigenspace identified by the same decomposition.
\end{solution}

\begin{problem}[Pairwise differences]
Find the matrix of the quadratic form $\sum_{i,j=1}^n(x_i-x_j)^2$ and
all its eigenvalues and eigenspaces.
\end{problem}
\begin{solution}
Expansion gives
\[
\sum_{i,j}(x_i-x_j)^2
=2n\sum_i x_i^2-2\left(\sum_i x_i\right)^2
=x^\top(2nI-2\mathbf1\mathbf1^\top)x=2n x^\top Cx.
\]
For $n>1$, the eigenvalue on $\R\mathbf1$ is $0$ and on
$\mathbf1^\perp$ it is $2n$, with multiplicities $1$ and $n-1$.
Thus an orthonormal eigenbasis is the normalized constant vector and
any orthonormal basis of its perpendicular complement.
For $n=1$, the form is identically zero and its only eigenvalue is $0$.
The identity also shows that total pairwise squared differences equal
$2n\sum_i(x_i-\bar x)^2$.
\end{solution}

\subsection*{Variational formulas and compression}
In the next ten exercises, $A$ is real symmetric of order $n$,
$\lambda_1\geqslant\cdots\geqslant\lambda_n$ are its eigenvalues,
$s_1,\ldots,s_n$ is a corresponding orthonormal eigenbasis, and every
Rayleigh quotient is evaluated only at a nonzero vector.

\begin{problem}[Excluding selected eigenvectors]\label{ex:12:exclude}
Let $R_j=(s_1\ \cdots\ s_j)$ and $T_j=(s_j\ \cdots\ s_n)$.
Prove
\[
\lambda_k=\max_{R_{k-1}^\top x=0}\rho_A(x)
=\min_{T_{k+1}^\top x=0}\rho_A(x).
\]
An empty matrix of constraints imposes no restriction.
\end{problem}
\begin{solution}
Write $x=\sum_i c_is_i$. In the first feasible set,
$c_1=\cdots=c_{k-1}=0$, so
$\rho_A(x)=\sum_{i=k}^n\lambda_i c_i^2/\sum_{i=k}^n c_i^2
\leqslant\lambda_k$. Equality occurs at $x=s_k$.
In the second feasible set $c_{k+1}=\cdots=c_n=0$, so the analogous
weighted average of $\lambda_1,\ldots,\lambda_k$ is at least
$\lambda_k$, again attained at $s_k$.
For $k=1$ in the maximum and $k=n$ in the minimum, the corresponding
coefficient restriction is empty, giving the usual extreme Rayleigh bounds.
\end{solution}

\begin{problem}[Arbitrary linear restrictions]
For every $B\in\R^{n\times(k-1)}$ and $C\in\R^{n\times(n-k)}$, prove
\[
\min_{C^\top x=0}\rho_A(x)\leqslant\lambda_k
\leqslant\max_{B^\top x=0}\rho_A(x).
\]
Do not assume that the columns of $B$ or $C$ are independent.
\end{problem}
\begin{solution}
Let $R=(s_1\ \cdots\ s_k)$. The matrix $B^\top R$ has $k$ columns
and at most $k-1$ rows, so some nonzero $u$ satisfies $B^\top Ru=0$.
Then $x=Ru\ne0$ is feasible and
$\rho_A(x)=\sum_{i=1}^k\lambda_i u_i^2/\sum_i u_i^2\geqslant\lambda_k$.
This proves the right inequality. For the left, take
$T=(s_k\ \cdots\ s_n)$, whose $n-k+1$ columns exceed the row count
of $C^\top T$. Choose $v\ne0$ in its kernel. Then $Tv\ne0$ is feasible
and has Rayleigh quotient at most $\lambda_k$. Dependence among the
constraint columns only increases these kernels, so it causes no problem.
\end{solution}

\begin{problem}[Fischer's min--max formulas]
Prove
\[
\lambda_k=\min_{B^\top B=I_{k-1}}\ \max_{B^\top x=0}\rho_A(x)
=\max_{C^\top C=I_{n-k}}\ \min_{C^\top x=0}\rho_A(x),
\]
where $B$ and $C$ have the dimensions in the preceding problem.
\end{problem}
\begin{solution}
For every admissible $B$, the preceding problem makes the inner maximum
at least $\lambda_k$. Choosing $B=(s_1\ \cdots\ s_{k-1})$ makes it
exactly $\lambda_k$ by Exercise~\ref{ex:12:exclude}. Thus the outer
minimum is $\lambda_k$ and is attained.
For every admissible $C$, the inner minimum is at most $\lambda_k$.
Choosing $C=(s_{k+1}\ \cdots\ s_n)$ attains equality, so the outer
maximum is $\lambda_k$. The empty-column cases work by the stated
constraint convention.
\end{solution}

\begin{problem}[Eigenvalue monotonicity]
If $B$ is positive semidefinite, prove
$\lambda_k(A+B)\geqslant\lambda_k(A)$ for every $k$.
Show the inequalities are strict when $B$ is positive definite.
\end{problem}
\begin{solution}
For every nonzero $x$,
$\rho_{A+B}(x)=\rho_A(x)+\rho_B(x)
\geqslant\rho_A(x)+\lambda_n(B)$.
Apply the maximum over $k$-dimensional subspaces of the minimum Rayleigh
quotient on each subspace to obtain
$\lambda_k(A+B)\geqslant\lambda_k(A)+\lambda_n(B)$.
For $B\succeq0$, $\lambda_n(B)\geqslant0$; for $B\succ0$, it is
strictly positive. Thus positive definiteness gives strict increase for
every ordered eigenvalue, even if $A$ itself is indefinite.
\end{solution}

\begin{problem}[$*$ Poincar\'e separation]\label{ex:12:separation}
Let $G\in\R^{n\times r}$ satisfy $G^\top G=I_r$. If
$\mu_1\geqslant\cdots\geqslant\mu_r$ are the eigenvalues of $G^\top AG$,
prove $\lambda_{n-r+k}\leqslant\mu_k\leqslant\lambda_k$.
\end{problem}
\begin{solution}
The map $y\mapsto Gy$ preserves lengths, dimensions and Rayleigh
quotients. By max--min, $\mu_k$ is the maximum, over $k$-dimensional
subspaces contained in $\operatorname{col}(G)$, of the minimum of
$\rho_A$ on that subspace. Removing the restriction that the subspace
lie in $\operatorname{col}(G)$ can only increase the maximum, giving
$\mu_k\leqslant\lambda_k$.
By min--max, $\mu_k$ is the minimum of the maximum over
$(r-k+1)$-dimensional subspaces in $\operatorname{col}(G)$.
Allowing all subspaces of that dimension can only decrease the minimum,
giving the bound $\mu_k\geqslant\lambda_{n-r+k}$.
Both arguments include $r=n$ and the endpoint eigenvalues.
\end{solution}

\begin{problem}[Compression by an orthogonal projector]
Let $P=P^\top=P^2$ have rank $r\geqslant1$. Write the eigenvalues of
$PAP$, apart from $n-r$ zero eigenvalues, as
$\mu_1\geqslant\cdots\geqslant\mu_r$. Prove
$\lambda_{n-r+k}\leqslant\mu_k\leqslant\lambda_k$.
\end{problem}
\begin{solution}
Choose an orthonormal basis $G$ of the range of $P$ and complete it by
$F$ to an orthogonal matrix. Then $P=GG^\top$ and
\[
(G\ F)^\top PAP(G\ F)
=\begin{pmatrix}G^\top AG&0\\0&0\end{pmatrix}.
\]
Thus the $r$ listed eigenvalues are those of $G^\top AG$, including
any zeros among those $r$ values. Apply the preceding separation theorem.
The wording does not require $PAP$ to have rank exactly $r$.
\end{solution}

\begin{problem}[Principal submatrices]
For any $r\times r$ principal submatrix $A_J$ of $A$, prove
$\lambda_{n-r+k}(A)\leqslant\lambda_k(A_J)\leqslant\lambda_k(A)$.
Write the full interlacing chain when $r=n-1$.
\end{problem}
\begin{solution}
Let $G$ consist of the coordinate columns indexed by $J$.
Then $G^\top G=I_r$ and $G^\top AG=A_J$, so separation applies.
For $r=n-1$, putting $\mu_k=\lambda_k(A_J)$ gives
$\lambda_{k+1}\leqslant\mu_k\leqslant\lambda_k$, or in full order,
\[
\lambda_n\leqslant\mu_{n-1}\leqslant\lambda_{n-1}
\leqslant\mu_{n-2}\leqslant\cdots\leqslant\mu_1\leqslant\lambda_1.
\]
The selected indices need not be the first $r$ indices.
\end{solution}

\begin{problem}[Trace bounds and attaining subspaces]\label{ex:12:trace}
Prove that the maximum and minimum of $\tr(Q^\top AQ)$ over
$Q^\top Q=I_r$ are respectively $\sum_{k=1}^r\lambda_k$ and
$\sum_{k=1}^r\lambda_{n-r+k}$. Deduce the same bounds for the trace
of the leading $r\times r$ principal submatrix.
\end{problem}
\begin{solution}
Let $\mu_1,\ldots,\mu_r$ be the eigenvalues of $Q^\top AQ$.
Separation gives $\lambda_{n-r+k}\leqslant\mu_k\leqslant\lambda_k$.
Summing yields the two bounds since $\tr(Q^\top AQ)=\sum_k\mu_k$.
The upper bound is attained by $Q=(s_1\ \cdots\ s_r)$; the lower by
$Q=(s_{n-r+1}\ \cdots\ s_n)$.
Taking $Q=(e_1\ \cdots\ e_r)$ gives the principal-submatrix trace.
This proof works for indefinite $A$ as well as for covariance matrices.
\end{solution}

\begin{problem}[Determinant bounds for compressed matrices]
Now assume $A\succ0$. Find the maximum and minimum of
$\det(Q^\top AQ)$ under $Q^\top Q=I_r$, and deduce bounds for the
determinant of the leading $r\times r$ principal submatrix.
\end{problem}
\begin{solution}
The compressed matrix is positive definite, since
$y^\top Q^\top AQy>0$ for $y\ne0$.
Its eigenvalues satisfy the positive bounds
$0<\lambda_{n-r+k}\leqslant\mu_k\leqslant\lambda_k$.
Multiplication, which preserves inequalities between positive numbers,
gives
\[
\prod_{k=1}^r\lambda_{n-r+k}\leqslant\det(Q^\top AQ)
\leqslant\prod_{k=1}^r\lambda_k.
\]
Both bounds are attained by bottom or top eigenvectors, respectively.
The principal-submatrix result follows by choosing the first $r$
coordinate vectors as columns of $Q$.
\end{solution}

\begin{problem}[A diagonal-product consequence]
For $A=(a_{ij})\succ0$, prove
$\prod_{k=1}^r\lambda_{n-r+k}\leqslant\prod_{k=1}^r a_{kk}$
for $1\leqslant r\leqslant n$.
\end{problem}
\begin{solution}
Let $B$ be the leading principal submatrix. The preceding determinant
bound gives $\prod_{k=1}^r\lambda_{n-r+k}\leqslant\det B$.
Write $B=R^\top R$ by applying Gram--Schmidt to a square root of $B$,
with $R$ upper triangular. Then
$\det B=\prod_{j=1}^r r_{jj}^2$, whereas
$b_{jj}=\sum_{i\leqslant j}r_{ij}^2\geqslant r_{jj}^2$.
Thus $\det B\leqslant\prod_j b_{jj}=\prod_{j=1}^r a_{jj}$.
Combining the two inequalities proves the claim, including $r=n$.
\end{solution}

\subsection*{SVD structure and PCA calculations}
\begin{problem}[Constructing the right singular directions]
Let $A\in\R^{m\times n}$ have rank $r$, and let $P=(P_1\ P_2)$ be
orthogonal with
$P^\top AA^\top P=\operatorname{diag}(D^2,0)$, where $D$ is an
invertible diagonal $r\times r$ matrix. Put $Q_1=A^\top P_1D^{-1}$
and let $Q_2\in\R^{n\times(n-r)}$ satisfy $Q_1^\top Q_2=0$.
Prove $P^\top A(Q_1\ Q_2)=\begin{pmatrix}D&0\\0&0\end{pmatrix}$.
\end{problem}
\begin{solution}
$Q_1^\top Q_1=D^{-1}P_1^\top AA^\top P_1D^{-1}=I_r$.
Also $P_2^\top AA^\top P_2=0$, so
$\norm{A^\top P_2}_F^2=0$ and $P_2^\top A=0$.
The relation $AA^\top P_1=P_1D^2$ gives
$AQ_1=P_1D$. In particular $P_1^\top AQ_1=D$.
Since $Q_1^\top=D^{-1}P_1^\top A$, the assumed
$Q_1^\top Q_2=0$ gives $P_1^\top AQ_2=0$.
The other two blocks vanish by $P_2^\top A=0$.
The requested block identity follows; it does not require $Q_2$ to
be orthonormal. To obtain an orthogonal SVD matrix $Q$, choose $Q_2$
to be an orthonormal basis of $\ker Q_1^\top$.
\end{solution}

\begin{problem}[Grouping equal singular values]
Suppose $A=\sum_{i=1}^r\sigma_i u_iv_i^\top$ is an SVD.
For each distinct positive singular value $a_j$, let
$W_j=\sum_{\{i:\sigma_i=a_j\}}u_iv_i^\top$.
Prove $W_jW_j^\top W_j=W_j$, and for $j\ne l$,
$W_l^\top W_j=0$ and $W_lW_j^\top=0$.
\end{problem}
\begin{solution}
Let $I_j$ be the index set in the definition. Orthonormality gives
\[
W_j^\top W_j=\sum_{i\in I_j}v_iv_i^\top,
\qquad W_jW_j^\top=\sum_{i\in I_j}u_iu_i^\top.
\]
Multiplying the second equality by $W_j$ makes all cross terms zero
and keeps $u_iv_i^\top$ for each $i\in I_j$, giving $W_j$.
For distinct groups the index sets are disjoint. In
$W_l^\top W_j$, every summand contains $u_h^\top u_i=0$ with
$h\in I_l,i\in I_j$; in $W_lW_j^\top$, every summand instead contains
$v_h^\top v_i=0$. Thus both products vanish.
\end{solution}

\begin{problem}[Range and kernel in an SVD]
For $A=P_1DQ_1^\top$ with $D$ invertible and $P=(P_1\ P_2)$,
$Q=(Q_1\ Q_2)$ orthogonal, prove
$\operatorname{col}(A)=\operatorname{col}(P_1)$ and
$\ker A=\operatorname{col}(Q_2)$.
\end{problem}
\begin{solution}
Every output of $A$ is in $\operatorname{col}(P_1)$, while
$AQ_1D^{-1}=P_1$, giving the reverse inclusion.
Also $Ax=0$ holds exactly when $Q_1^\top x=0$, because $P_1$ has
independent columns and $D$ is invertible. The orthogonal expansion
$x=Q_1Q_1^\top x+Q_2Q_2^\top x$ then shows that this kernel is exactly
$\operatorname{col}(Q_2)$. Thus discarded zero singular directions
describe a null space, not missing positive-variance directions.
\end{solution}

\begin{problem}[Two-variable PCA from observations]\label{ex:12:numerical-pca}
For $D=\begin{pmatrix}12&21\\11&22\\9&18\\8&19\end{pmatrix}$, compute
the mean, centered matrix, covariance, principal directions, score
columns and explained-variance proportions. Reconstruct the observations
with one component and compute the total squared reconstruction error.
\end{problem}
\begin{solution}
The mean is $(10,20)^\top$ and
$Z=\begin{pmatrix}2&1\\1&2\\-1&-2\\-2&-1\end{pmatrix}$.
Thus $S=\frac13\begin{pmatrix}10&8\\8&10\end{pmatrix}$.
The sum and difference directions are
$v_1=(1,1)^\top/\sqrt2$, $v_2=(1,-1)^\top/\sqrt2$, with eigenvalues
$6,2/3$. Their scores are $(3,3,-3,-3)^\top/\sqrt2$ and
$(1,-1,1,-1)^\top/\sqrt2$, respectively.
The proportions are $9/10$ and $1/10$.
Keeping the first component gives
\[
D_1=\begin{pmatrix}23/2&43/2\\23/2&43/2\\17/2&37/2\\17/2&37/2\end{pmatrix}.
\]
Subtracting from $D$ gives residual rows
$(1/2,-1/2),(-1/2,1/2),(1/2,-1/2),(-1/2,1/2)$.
Their squared lengths sum to $2$, equal to $(n-1)\lambda_2=3(2/3)$.
The best-rank theorem proves that no other rank-one centered fit has
smaller squared error.
\end{solution}

\begin{problem}[More variables than observations]
For $D=\begin{pmatrix}1&2&3&4\\2&4&6&8\\3&6&9&12\end{pmatrix}$,
compute PCA using the observation Gram matrix, then recover the loading
and scores. State the covariance rank and the number of nonzero components.
\end{problem}
\begin{solution}
Put $a=(1,2,3,4)^\top$ and $u=(-1,0,1)^\top$.
The mean is $2a$ and $Z=ua^\top$, where $a^\top a=30$, $u^\top u=2$.
Thus $G=ZZ^\top/2=15uu^\top$ has unit eigenvector $u/\sqrt2$ with
eigenvalue $30$, and has two zero eigenvalues.
The recovered loading is
\[
v=\frac{Z^\top(u/\sqrt2)}{\sqrt{2\cdot30}}
=\frac{a\sqrt2}{\sqrt{60}}=\frac a{\sqrt{30}}.
\]
Its scores are $Zv=\sqrt{30}\,u$.
The covariance is $S=aa^\top$, of rank one with eigenvalues $30,0,0,0$.
One component explains all variation and reconstructs $Z$ exactly.
The other loadings may be any orthonormal basis of $a^\perp$ and have
identically zero scores.
\end{solution}

\begin{problem}[Covariance and correlation PCA]\label{ex:12:standardize}
For the covariance matrix $S=\begin{pmatrix}9&3\\3&1\end{pmatrix}$,
compute covariance PCA and correlation PCA. Express the first standardized
score as a linear combination of the original centered variables.
\end{problem}
\begin{solution}
$S=aa^\top$ with $a=(3,1)^\top$.
Covariance PCA has eigenvalues $10,0$ and first loading
$(3,1)^\top/\sqrt{10}$.
The standard-deviation matrix is $D_s=\operatorname{diag}(3,1)$.
Thus $R=D_s^{-1}SD_s^{-1}=\begin{pmatrix}1&1\\1&1\end{pmatrix}$,
with eigenvalues $2,0$ and first loading $(1,1)^\top/\sqrt2$ in
standardized coordinates. For a centered observation $(z_1,z_2)$,
its standardized score is $(z_1/3+z_2)/\sqrt2$.
Both methods retain all variation with one component, but their unit
loading directions and score scales are different. Standardization is
valid because both marginal variances are positive, despite singularity
of $S$.
\end{solution}

\begin{problem}[A tie at the cutoff]
For $S=\operatorname{diag}(4,4,1)$, describe all optimal one-dimensional
and two-dimensional PCA subspaces and their explained-variance proportions.
For $Z=\operatorname{diag}(2,2,1)$, describe all best rank-one Frobenius
approximations and determine the minimum squared error.
\end{problem}
\begin{solution}
The total variance is $9$. Every line in $\operatorname{span}(e_1,e_2)$
is an optimal first direction, explaining $4/9$. The optimal
two-dimensional space is uniquely $\operatorname{span}(e_1,e_2)$,
explaining $8/9$, because the gap after the second eigenvalue is positive.
An orthonormal basis within that plane can still be rotated arbitrarily.
For the separate approximation problem, $Z^\top Z=S$.
The low-rank proof shows that an optimal rank-one matrix must equal
$ZP$, where $P=vv^\top$ projects onto an optimal line with
$v=(a,b,0)^\top$ and $a^2+b^2=1$.
Thus every best approximation is $2vv^\top$ and has squared error
$4+1=5$. Conversely each such choice attains that error.
A rank-zero competitor has error $9$, so no additional zero-rank
minimizer is missing.
\end{solution}

\begin{problem}[The affine reconstruction problem]
For observations $d_1,\ldots,d_n\in\R^p$, fix a direction subspace $W$
with projector $P$. Prove that an affine translate $a+W$ minimizes the
sum of squared distances for this direction precisely when it contains
the observation mean. Then identify the best direction space of dimension
$k$ in terms of the sample covariance.
\end{problem}
\begin{solution}
Write $d_i-a=(d_i-\bar d)+(\bar d-a)$.
Since $\sum_i(d_i-\bar d)=0$, expansion gives
\[
\sum_i\norm{(I-P)(d_i-a)}^2
=\sum_i\norm{(I-P)(d_i-\bar d)}^2
+n\norm{(I-P)(\bar d-a)}^2.
\]
The first term does not depend on $a$. The second is zero exactly when
$\bar d-a\in W$, equivalently $\bar d\in a+W$.
For such a translate, the objective is
$(n-1)[\tr S-\tr(PS)]$.
The trace variational theorem makes it minimal for a span of the top
$k$ eigenvectors, with the cutoff-tie choices described there.
The minimum is $(n-1)\sum_{i>k}\lambda_i(S)$.
\end{solution}

\begin{problem}[Equicorrelation interpreted as PCA]
For the $p\times p$ equicorrelation covariance with $p\geqslant2$ and
$-1/(p-1)<t<1$, determine the leading PCA subspace and its variance
proportion when $t>0$, $t=0$, and $t<0$. What happens at the two
positive-semidefinite endpoints?
\end{problem}
\begin{solution}
The total variance is $p$. The constant direction has eigenvalue
$a=1+(p-1)t$, and its perpendicular complement has eigenvalue $b=1-t$.
If $t>0$, $a>b$ and the leading subspace is uniquely
$\R\mathbf1$, explaining $a/p$.
If $t=0$, all eigenvalues equal $1$, so every line is optimal and explains
$1/p$. If $t<0$, $b>a$, and any line in $\mathbf1^\perp$ is optimal,
explaining $b/p$; the entire leading eigenspace has dimension $p-1$.
At $t=1$, $a=p,b=0$, so one constant-direction component captures
everything. At $t=-1/(p-1)$, $a=0,b=p/(p-1)$, and the $p-1$
perpendicular directions together capture all variance, each contributing
$1/(p-1)$. For $p=2$ the perpendicular eigenspace is one-dimensional.
\end{solution}

\needspace{10\baselineskip}
\subsection*{Further matrix inequalities}
In this supplement all matrices are real. For symmetric matrices,
\(A\geqslant B\) means \(A-B\) is positive semidefinite, and
\(A>B\) means \(A-B\) is positive definite. Orders are positive integers.
\(\det\) denotes determinant; \(\norm{\cdot}_F\) denotes the Frobenius norm.

\needspace{8\baselineskip}
\subsection*{Schwarz, trace, and rank}

\begin{problem}\label{ex:12:ineq-ma-01}
Prove \((a^\top b)^2\leqslant(a^\top a)(b^\top b)\), with equality
exactly when \(a,b\) are dependent, by each of these methods:
(a) use \(ab^\top-ba^\top\); (b) for \(b\ne0\), use
\(I-bb^\top/(b^\top b)\); (c) use the positive semidefinite matrices
\(c_ic_i^\top\), where \(c_i=(a_i,b_i)^\top\).
\end{problem}
\begin{solution}
(a) Expanding the sum of squared entries gives
\[
\norm{ab^\top-ba^\top}_F^2
=\sum_{i,j}(a_i b_j-b_i a_j)^2
=2\big(\norm{a}^2\norm{b}^2-(a^\top b)^2\big)\geqslant0.
\]
Equality means \(ab^\top=ba^\top\). If \(b\ne0\), multiply by \(b\)
to get \(a\norm{b}^2=b(a^\top b)\), so the vectors are dependent;
if \(b=0\), dependence and equality both hold. The converse is immediate.
(b) The indicated matrix \(P\) is symmetric idempotent, so
\(a^\top Pa=\norm{Pa}^2\geqslant0\). Expanding gives the inequality,
and equality means \(Pa=0\), exactly \(a\in\operatorname{span}(b)\).
(c) The sum is
\(\sum_i c_ic_i^\top=\begin{pmatrix}\norm{a}^2&a^\top b\\a^\top b&\norm{b}^2\end{pmatrix}\),
a positive semidefinite Gram matrix, whose determinant is nonnegative.
It is singular exactly when the two columns of \((a\mid b)\) are
dependent: \(\ker(C^\top C)=\ker C\) proves this precisely.
Thus equality is again equivalent to dependence.
\end{solution}

\begin{problem}\label{ex:12:ineq-ma-02}
Let \(A=(a_{ij})\) be positive semidefinite. Prove
\((x^\top Ay)^2\leqslant(x^\top Ax)(y^\top Ay)\), with equality
exactly when \(Ax,Ay\) are dependent. Deduce
\(|a_{ij}|\leqslant\max_k a_{kk}\).
\end{problem}
\begin{solution}
Apply Schwarz to \(A^{1/2}x,A^{1/2}y\).
Equality means that a nontrivial combination \(z=ax+by\) belongs to
\(\ker A^{1/2}=\ker A\), exactly the dependence of \(Ax,Ay\).
With \(x=e_i,y=e_j\), obtain \(|a_{ij}|\leqslant\sqrt{a_{ii}a_{jj}}\).
The diagonal entries are nonnegative, and their geometric mean is no
larger than their maximum, proving the result.
\end{solution}

\begin{problem}\label{ex:12:ineq-ma-03}
For positive definite \(A\), prove
\((x^\top y)^2\leqslant(x^\top A^{-1}x)(y^\top Ay)\), with equality
exactly when \(A^{-1}x,y\) are dependent.
For fixed \(x\ne0\), show
\[
\frac1{x^\top A^{-1}x}
=\min_{y:\,y^\top x\ne0}\frac{y^\top Ay}{(y^\top x)^2}.
\]
Deduce, for positive definite \(A,B\),
\[
x^\top(A+B)^{-1}x\leqslant
\frac{(x^\top A^{-1}x)(x^\top B^{-1}x)}{x^\top(A^{-1}+B^{-1})x}.
\]
\end{problem}
\begin{solution}
Schwarz applied to \(A^{-1/2}x,A^{1/2}y\) gives the first inequality;
its equality condition becomes dependence after multiplication by
\(A^{-1/2}\). Dividing by the positive quantities gives the proposed
lower bound, attained at \(y=A^{-1}x\), for which \(y^\top x>0\).
Let \(\psi(A)=1/(x^\top A^{-1}x)\). For each admissible \(y\),
\[
\frac{y^\top(A+B)y}{(y^\top x)^2}
\geqslant\psi(A)+\psi(B).
\]
Taking the minimum gives \(\psi(A+B)\geqslant\psi(A)+\psi(B)\).
All terms are positive, so taking reciprocals and combining the two
fractions proves the stated bound.
\end{solution}

\begin{problem}\label{ex:12:ineq-ma-04}
(a) Prove \((\sum_i x_i)^2\leqslant n\sum_i x_i^2\), with equality
exactly when all \(x_i\) agree.
(b) If all eigenvalues of a real square \(A\) are real, prove
\[|\tr A/n|\leqslant\sqrt{\tr(A^2)/n},\]
with equality exactly when all eigenvalues agree.
(c) If \(A\ne0\) is symmetric, show
\(\operatorname{rank}A\geqslant(\tr A)^2/\tr(A^2)\).
\end{problem}
\begin{solution}
(a) Apply Schwarz to \((x_i)\) and the all-ones vector; dependence means
all coordinates agree.
(b) The trace identities \(\tr A=\sum_i\lambda_i\) and
\(\tr(A^2)=\sum_i\lambda_i^2\) reduce the claim to (a).
For completeness these identities require no diagonalizability: choose
one real eigenvector as the first basis vector, obtaining a block upper
triangular matrix \(\begin{pmatrix}\lambda&*\\0&C\end{pmatrix}\).
Its characteristic polynomial factors, and its trace and squared trace
contributions are \(\lambda+\tr C\) and \(\lambda^2+\tr(C^2)\).
Induction proves both identities because the remaining eigenvalues stay real.
(c) Apply (a) to the \(r=\operatorname{rank}A\) nonzero eigenvalues:
\((\tr A)^2\leqslant r\tr(A^2)\). The denominator is positive because
\(A\ne0\) is symmetric.
\end{solution}

\begin{problem}\label{ex:12:ineq-ma-05}
For real matrices \(A,B\) of the same size, prove
\begin{enumerate}
\item \((\tr(A^\top B))^2\leqslant\tr(A^\top A)\tr(B^\top B)\),
with equality exactly when one matrix is a scalar multiple of the other;
\item \(\tr((A^\top B)^2)\leqslant\tr(A^\top AB^\top B)\),
with equality exactly when \(AB^\top\) is symmetric;
\item \(\tr((A^\top B)^2)\leqslant\tr(AA^\top BB^\top)\),
with equality exactly when \(A^\top B\) is symmetric.
\end{enumerate}
\end{problem}
\begin{solution}
(a) This is Schwarz for the Frobenius inner product; viewing all entries
as one vector gives the stated equality condition, including zero matrices.
For any real square \(M\), direct expansion gives
\[
\norm{M-M^\top}_F^2=2\big(\tr(M^\top M)-\tr(M^2)\big).
\]
(b) Apply this to \(M=AB^\top\). Cyclicity gives
\(\tr(M^2)=\tr((A^\top B)^2)\) (also use invariance under transpose),
and \(\tr(M^\top M)=\tr(A^\top AB^\top B)\).
The nonnegative difference vanishes exactly when \(M=M^\top\).
(c) Apply the same identity to \(M=A^\top B\), for which
\(\tr(M^\top M)=\tr(AA^\top BB^\top)\).
\end{solution}

\begin{problem}\label{ex:12:ineq-ma-06}
For any real square \(A\), prove \(\tr(A^2)\leqslant\tr(A^\top A)\),
with equality exactly when \(A\) is symmetric.
\end{problem}
\begin{solution}
Expanding the Frobenius square yields
\(\norm{A-A^\top}_F^2=2\tr(A^\top A)-2\tr(A^2)\).
It is nonnegative and is zero exactly when \(A=A^\top\).
\end{solution}

\needspace{8\baselineskip}
\subsection*{Determinants and positive order}

\begin{problem}\label{ex:12:ineq-ma-07}
For positive semidefinite \(n\times n\) matrices \(A,B\), prove
\(\det(A+B)\geqslant\det A+\det B\). Determine every equality case.
\end{problem}
\begin{solution}
If both are singular, the right side is zero and the inequality follows
from positive semidefiniteness of their sum. Equality then means the sum
is singular. Otherwise suppose \(A\) is positive definite, exchanging
the names if needed. Put \(C=A^{-1/2}BA^{-1/2}\), with eigenvalues
\(\lambda_i\geqslant0\). Dividing the proposed inequality by \(\det A>0\)
reduces it to
\(\prod_i(1+\lambda_i)\geqslant1+\prod_i\lambda_i\),
which follows by expanding the product into nonnegative terms.
For \(n=1\) there is always equality. For \(n\geqslant2\), the omitted
terms include \(\sum_i\lambda_i\), so equality in this nonsingular case
requires every \(\lambda_i=0\), or \(B=0\).
Thus all equality cases are: \(n=1\), or \(A+B\) singular, or \(A=0\),
or \(B=0\).
\end{solution}

\begin{problem}\label{ex:12:ineq-ma-08}
If \(A\geqslant B\geqslant0\), show \(\det A\geqslant\det B\).
Prove that equality holds exactly when both matrices are singular, or
when both are nonsingular and \(A=B\).
\end{problem}
\begin{solution}
If \(B\) is singular, its determinant is zero and equality holds exactly
when \(A\) is singular too. If \(B\) is positive definite, write
\(A=B^{1/2}(I+C)B^{1/2}\), where
\(C=B^{-1/2}(A-B)B^{-1/2}\geqslant0\).
Then \(\det A/\det B=\prod_i(1+\lambda_i(C))\geqslant1\),
with equality precisely when all eigenvalues of \(C\) vanish, or \(A=B\).
\end{solution}

\begin{problem}\label{ex:12:ineq-ma-09}
If \(0\leqslant A\leqslant I_n\), prove \(A=I_n\iff\det A=1\).
\end{problem}
\begin{solution}
The eigenvalues lie in \([0,1]\), since for a unit eigenvector \(v\),
\(0\leqslant v^\top Av\leqslant1\). Their product is one exactly
when every eigenvalue is one. Orthogonal diagonalization then gives
\(A=I_n\). The converse is immediate.
\end{solution}

\begin{problem}\label{ex:12:ineq-ma-11}
Prove the arithmetic–geometric mean inequality
\((\prod_i\lambda_i)^{1/n}\leqslant\sum_i\lambda_i/n\) for nonnegative
\(\lambda_i\), with equality exactly when all are equal.
Deduce \((\det A)^{1/n}\leqslant\tr A/n\) for positive semidefinite
\(A\), and prove that equality holds exactly when \(A=\alpha I_n\), \(\alpha\geqslant0\).
\end{problem}
\begin{solution}
For \(t>0\), the function \(t-1-\log t\) has derivative \(1-1/t\),
so its unique minimum is zero at \(t=1\). If every \(\lambda_i>0\),
let \(a=\sum_i\lambda_i/n\). Sum
\(\log(\lambda_i/a)\leqslant\lambda_i/a-1\) to obtain
\(\log(\prod_i\lambda_i/a^n)\leqslant0\).
Equality requires every \(\lambda_i=a\). If some \(\lambda_i=0\), the
product is zero, and equality requires that their nonnegative sum be zero,
so all are zero. Apply this scalar result to the eigenvalues of \(A\).
Equality makes its orthogonal diagonal form \(\alpha I_n\), and hence
\(A=\alpha I_n\).
\end{solution}

\begin{problem}\label{ex:12:ineq-ma-12}
For positive semidefinite \(A\) and positive definite \(X\) with
\(\det X=1\), prove \(\tr(AX)/n\geqslant(\det A)^{1/n}\).
Show that equality holds exactly when \(A=0\), or when \(A\) is positive
definite and \(X=(\det A)^{1/n}A^{-1}\).
\end{problem}
\begin{solution}
Apply the preceding arithmetic–geometric mean inequality to
\(C=X^{1/2}AX^{1/2}\). Its trace is \(\tr(AX)\), and its determinant
is \(\det A\). Equality means \(C=\alpha I_n\).
If \(\alpha=0\), invertibility of \(X^{1/2}\) gives \(A=0\).
Otherwise \(A=\alpha X^{-1}>0\); taking determinants yields
\(\alpha=(\det A)^{1/n}\), hence the stated \(X\).
Each of these choices attains equality by substitution.
\end{solution}

\begin{problem}\label{ex:12:ineq-ma-13}
For nonzero positive semidefinite \(A,B\), prove
\[
(\det(A+B))^{1/n}\geqslant(\det A)^{1/n}+(\det B)^{1/n}.
\]
Determine every equality case using the preceding trace representation.
\end{problem}
\begin{solution}
If \(S=A+B\) is singular, a nonzero \(z\in\ker S\) satisfies
\(z^\top Az+z^\top Bz=0\), so both terms vanish and both \(A,B\)
are singular. Equality is then zero equals zero.
If \(S>0\), choose \(X=(\det S)^{1/n}S^{-1}\), with \(\det X=1\).
The preceding exercise gives
\[
(\det S)^{1/n}=\tr(SX)/n=\tr(AX)/n+\tr(BX)/n
\geqslant(\det A)^{1/n}+(\det B)^{1/n}.
\]
Because neither summand matrix is zero, equality requires both to be
positive definite and the same \(X\) to be their equality choice.
Thus \((\det A)^{1/n}A^{-1}=(\det B)^{1/n}B^{-1}\), equivalent to
\(A=cB\) for some \(c>0\). Such proportional matrices attain equality.
Altogether equality holds exactly when \(A+B\) is singular or \(A=cB\)
with \(c>0\).
\end{solution}

\begin{problem}\label{ex:12:ineq-ma-14}
For positive semidefinite \(A,B\), prove
\(0\leqslant\tr(AB)\leqslant\tr A\,\tr B\) and
\[\sqrt{\tr(AB)}\leqslant\tfrac12(\tr A+\tr B).\]
Determine every equality case in the second inequality.
\end{problem}
\begin{solution}
Diagonalize \(A\) orthogonally to \(\operatorname{diag}(\lambda_i)\),
and write \(C\) for \(B\) in the same coordinates. Then
\(\tr(AB)=\sum_i\lambda_i c_{ii}\geqslant0\), and
\[
\tr A\,\tr B-\tr(AB)=\sum_{i\ne j}\lambda_i c_{jj}\geqslant0.
\]
Also \(\sqrt{ab}\leqslant(a+b)/2\) for nonnegative \(a,b\), since
\((\sqrt a-\sqrt b)^2\geqslant0\). Equality requires equal traces and
all the displayed cross terms zero. If either matrix is zero, equal
traces force both to be zero. Otherwise choose \(\lambda_k>0\).
Then \(c_{jj}=0\) for \(j\ne k\); the PSD entry bound forces the
corresponding rows and columns of \(C\) to vanish. As \(C\ne0\),
\(c_{kk}>0\), forcing \(\lambda_i=0\) for \(i\ne k\).
Equal traces now give \(\lambda_k=c_{kk}\).
Thus \(A=B=aa^\top\) for some nonzero vector \(a\); taking \(a=0\)
includes the zero case. These and only these choices attain equality.
\end{solution}

\begin{problem}\label{ex:12:ineq-ma-15}
For real \(m\times n\) matrices \(A,B\), prove
\[
(\det(A^\top B))^2\leqslant\det(A^\top A)\det(B^\top B).
\]
Prove equality holds exactly when one Gram matrix is singular, or
\(B=AQ\) for an invertible \(n\times n\) matrix \(Q\).
\end{problem}
\begin{solution}
If a Gram matrix is singular, its corresponding matrix has rank below
\(n\), so \(A^\top B\) is singular and both sides vanish.
Otherwise both have independent columns. Let
\(P=A(A^\top A)^{-1}A^\top\), the orthogonal projector onto
\(\operatorname{im}A\). Then \(0\leqslant B^\top PB\leqslant B^\top B\),
and determinant monotonicity gives
\[
\frac{(\det(A^\top B))^2}{\det(A^\top A)}
=\det(B^\top PB)\leqslant\det(B^\top B).
\]
Since the upper matrix is positive definite, equality requires
\(B^\top(I-P)B=0\). This is the Gram matrix of \((I-P)B\), so equality
means \(PB=B\). Thus \(B=AQ\); its independent columns force \(Q\)
to be invertible. Conversely this relation gives equality by determinant
multiplication.
\end{solution}

\needspace{8\baselineskip}
\subsection*{Inverse and power inequalities}

\begin{problem}\label{ex:12:ineq-ma-16}
For positive definite \(A,B\), prove
\(A>B\iff B^{-1}>A^{-1}\), where \(>\) means that the difference is
positive definite.
\end{problem}
\begin{solution}
If \(A>B\), then \(C=B^{-1/2}AB^{-1/2}>I\). Every eigenvalue of \(C\)
is greater than one, so \(I-C^{-1}>0\). Congruence by \(B^{-1/2}\)
gives \(B^{-1}-A^{-1}>0\). Apply the implication just proved to
\(B^{-1},A^{-1}\) to obtain the reverse direction.
\end{solution}

\begin{problem}\label{ex:12:ineq-ma-17}
(a) For \(a\leqslant t\leqslant b\), prove \(t^2-(a+b)t+ab\leqslant0\).
(b) Let \(A>0\) have order \(n\geqslant2\) and extreme eigenvalues
\(M\geqslant m>0\). Prove \((M+m)I-A-MmA^{-1}\geqslant0\), of rank
at most \(n-2\).
(c) Using the trace-square inequality and (b), prove for every unit \(x\)
\[
1\leqslant(x^\top Ax)(x^\top A^{-1}x)\leqslant\frac{(M+m)^2}{4Mm}.
\]
\end{problem}
\begin{solution}
(a) Factor the polynomial as \((t-a)(t-b)\leqslant0\).
(b) In an eigenbasis, the displayed matrix has entries
\(M+m-\lambda-Mm/\lambda\geqslant0\) by (a), with zero at each
extreme eigenvalue. If \(M>m\), these give at least two zero entries;
if \(M=m\), all entries vanish. The rank bound follows in both cases.
(c) Put \(D=A^{1/2}xx^\top A^{-1/2}\). Then
\(\tr(D^2)=1\) and
\(\tr(D^\top D)=(x^\top Ax)(x^\top A^{-1}x)\).
The trace-square inequality proves the lower bound.
Set \(u=x^\top Ax\), \(v=x^\top A^{-1}x\).
Part (b) gives \(u+Mm v\leqslant M+m\), whence
\[
Mmuv\leqslant(M+m)u-u^2
=\frac{(M+m)^2}{4}-\left(u-\frac{M+m}{2}\right)^2
\leqslant\frac{(M+m)^2}{4}.
\]
Division by \(Mm>0\) finishes the proof. For order one the two bounds
are both one, although the stated rank bound would need separate wording.
\end{solution}

\begin{problem}\label{ex:12:ineq-ma-19}
Suppose \(A\geqslant B\geqslant0\).
(a) Prove \(A^{1/2}\geqslant B^{1/2}\).
(b) Show that \(A^2\geqslant B^2\) need not hold.
(c) Prove that if \(AB=BA\), then \(A^k\geqslant B^k\) for every integer
\(k\geqslant2\).
\end{problem}
\begin{solution}
(a) Let \(C=A^{1/2}-B^{1/2}\), symmetric. Expanding the square gives
\(A-B=A^{1/2}C+CA^{1/2}-C^2\).
If \(Cx=\lambda x\), \(x\ne0\), and \(\lambda<0\), then
\[
x^\top(A-B)x=2\lambda x^\top A^{1/2}x-\lambda^2\norm{x}^2<0,
\]
contradicting \(A-B\geqslant0\). Thus every eigenvalue of \(C\) is
nonnegative, proving the claim, including singular matrices.
(b) Take \(A=\begin{pmatrix}2&1\\1&1\end{pmatrix}\),
\(B=\operatorname{diag}(1,0)\). The difference \(A-B\) is the all-ones
matrix and is positive semidefinite; \(A,B\) are positive semidefinite.
But \(A^2-B^2=\begin{pmatrix}4&3\\3&2\end{pmatrix}\) has value \(-2\)
on \((2,-3)^\top\), so it is not positive semidefinite.
(c) Each eigenspace of \(A\) is invariant under \(B\), because they
commute. Diagonalize the symmetric restriction of \(B\) within each
such eigenspace to obtain one orthonormal basis diagonalizing both.
Their diagonal entries satisfy \(\lambda_i\geqslant\mu_i\geqslant0\)
by the assumed order. Then \(\lambda_i^k\geqslant\mu_i^k\) proves the
matrix inequality after changing coordinates back.
\end{solution}

\needspace{8\baselineskip}
\subsection*{Logarithms and determinant concavity}

\begin{problem}\label{ex:12:ineq-ma-20}
For positive definite \(n\times n\) \(A\), prove
\(\log\det A\leqslant\tr A-n\), with equality exactly when \(A=I_n\).
\end{problem}
\begin{solution}
For each positive eigenvalue \(\lambda_i\), the scalar inequality
\(\log\lambda_i\leqslant\lambda_i-1\) was proved by differentiation
in the arithmetic–geometric mean exercise. Summing gives the stated bound.
Equality requires every \(\lambda_i=1\); orthogonal diagonalization
then gives \(A=I_n\), which also attains equality.
\end{solution}

\begin{problem}\label{ex:12:ineq-ma-21}
(a) For \(t>0\), \(0<\alpha<1\), prove
\(t^\alpha\leqslant\alpha t+1-\alpha\), with equality exactly at \(t=1\).
(b) For positive semidefinite \(A,B\) of the same size, prove
\[
(\det A)^\alpha(\det B)^{1-\alpha}
\leqslant\det(\alpha A+(1-\alpha)B).
\]
(c) Determine every equality case.
\end{problem}
\begin{solution}
(a) The derivative of \(\alpha t+1-\alpha-t^\alpha\) is
\(\alpha(1-t^{\alpha-1})\), negative before one and positive after one.
Its unique minimum is zero at one.
(b) If either matrix is singular, the left side is zero and the right
side nonnegative. If both are positive definite, set
\(C=B^{-1/2}AB^{-1/2}\), with positive eigenvalues \(\lambda_i\).
Divide both sides by \(\det B\); the claim becomes
\(\prod_i\lambda_i^\alpha\leqslant\prod_i(\alpha\lambda_i+1-\alpha)\),
which follows from (a).
(c) In the positive definite case, equality requires every \(\lambda_i=1\),
or \(A=B\). In the singular case equality means the mixture is singular.
For positive weights, its kernel is \(\ker A\cap\ker B\), since a zero
quadratic value in the sum forces both nonnegative terms to vanish.
Thus equality holds exactly when \(A=B\), or the two matrices have a
common nonzero kernel vector.
\end{solution}

\begin{problem}\label{ex:12:ineq-ma-22}
For positive definite matrices \(A_1,\ldots,A_k\) and positive weights
\(\alpha_i\) summing to one, prove
\[
\prod_{i=1}^k(\det A_i)^{\alpha_i}
\leqslant\det\left(\sum_{i=1}^k\alpha_iA_i\right).
\]
Determine when equality holds.
\end{problem}
\begin{solution}
Induct on \(k\). For \(k=1\) equality is automatic. For \(k>1\), let
\(s=\sum_{i<k}\alpha_i\) and \(C=\sum_{i<k}(\alpha_i/s)A_i\).
The induction hypothesis bounds the first \(k-1\) determinant factors
by \((\det C)^s\); the two-matrix inequality then gives
\((\det C)^s(\det A_k)^{1-s}\leqslant\det(sC+(1-s)A_k)\).
All intermediate matrices are positive definite. Equality requires
\(A_1=\cdots=A_{k-1}\) by induction and \(C=A_k\) by the two-matrix
equality case. Thus equality holds precisely when all \(A_i\) agree.
\end{solution}

\needspace{8\baselineskip}
\subsection*{Bordered matrices and volume bounds}

\begin{problem}\label{ex:12:ineq-ma-23}
Let \(A>0\) and \(Z=\begin{pmatrix}A&b\\b^\top&\alpha\end{pmatrix}\).
Prove (a) \(\det Z\leqslant\alpha\det A\), with equality exactly when
\(b=0\); (b) \(Z>0\iff\det Z>0\).
\end{problem}
\begin{solution}
With \(L=\begin{pmatrix}I&-A^{-1}b\\0&1\end{pmatrix}\), multiplication gives
\(L^\top ZL=\operatorname{diag}(A,\alpha-b^\top A^{-1}b)\), and \(\det L=1\).
Hence \(\det Z=\det A(\alpha-b^\top A^{-1}b)\).
Since \(A^{-1}>0\), the subtracted term is nonnegative and vanishes
exactly for \(b=0\), proving (a).
Congruence by invertible \(L\) preserves positive definiteness, so \(Z>0\)
exactly when \(\alpha-b^\top A^{-1}b>0\). As \(\det A>0\), this is
exactly \(\det Z>0\).
\end{solution}

\begin{problem}\label{ex:12:ineq-ma-24}
Let \(A\geqslant0\) and \(Z=\begin{pmatrix}A&b\\b^\top&\alpha\end{pmatrix}\).
Prove \(Z\geqslant0\) exactly when \(\alpha\geqslant0\) and
\(\alpha A-bb^\top\geqslant0\). Obtain the positive definite bordered
criterion and determinant bound of the preceding exercise as a special case.
\end{problem}
\begin{solution}
If \(Z\geqslant0\), then \(\alpha\geqslant0\) by its last diagonal entry.
If \(\alpha>0\), complete the square:
\[
\begin{pmatrix}x\\t\end{pmatrix}^{\!\top}Z\begin{pmatrix}x\\t\end{pmatrix}
=\alpha\left(t+\frac{b^\top x}{\alpha}\right)^2
+x^\top\left(A-\frac{bb^\top}{\alpha}\right)x.
\]
This proves both directions in this case. If \(\alpha=0\) and \(b\ne0\),
choose \(x=b\) and take \(t\) sufficiently negative to make the quadratic
expression negative. Thus \(Z\geqslant0\) forces \(b=0\); conversely
\(b=0\) makes \(Z=\operatorname{diag}(A,0)\geqslant0\).
The condition \(-bb^\top\geqslant0\) also forces \(b=0\), proving the
claimed equivalence at zero.
For \(A>0\), completing the square in \(x\) instead gives
\((x+tA^{-1}b)^\top A(x+tA^{-1}b)+(\alpha-b^\top A^{-1}b)t^2\).
It is positive definite exactly when the last coefficient is positive.
The same change of coordinates has determinant one, giving
\(\det Z=\det A(\alpha-b^\top A^{-1}b)\leqslant\alpha\det A\),
with equality exactly at \(b=0\), as required.
\end{solution}

\begin{problem}\label{ex:12:ineq-ma-25}
For \(A>0\), show that
\(Z=\begin{pmatrix}A&b\\b^\top&b^\top A^{-1}b\end{pmatrix}\)
is positive semidefinite and singular. Find a zero-eigenvalue eigenvector.
\end{problem}
\begin{solution}
Its quadratic form is
\((x+tA^{-1}b)^\top A(x+tA^{-1}b)\geqslant0\).
The nonzero vector \((-A^{-1}b,1)^\top\) is annihilated by both block
rows: the first gives \(-b+b=0\), the second
\(-b^\top A^{-1}b+b^\top A^{-1}b=0\).
Thus it is a zero-eigenvalue eigenvector and \(Z\) is singular.
\end{solution}

\begin{problem}\label{ex:12:ineq-ma-26}
(a) For positive definite \(A=(a_{ij})\), prove
\(\det A\leqslant\prod_i a_{ii}\), with equality exactly when \(A\)
is diagonal.
(b) Deduce for every real square \(A\) that
\((\det A)^2\leqslant\prod_i\sum_j a_{ij}^2\).
(c) Determine every equality case in (b).
\end{problem}
\begin{solution}
(a) Induct on order, the scalar case being equality. Partition
\(A=\begin{pmatrix}C&b\\b^\top&d\end{pmatrix}>0\), so \(C>0\).
The bordered determinant identity gives
\(\det A=\det C(d-b^\top C^{-1}b)\leqslant d\det C
\leqslant d\prod_{i<n}a_{ii}\).
Since all diagonal entries are positive, equality requires both \(b=0\)
and equality for \(C\), hence \(A\) diagonal. The converse follows directly.
(b) For nonsingular \(A\), apply (a) to \(AA^\top>0\), whose diagonal
entries are its row squared norms and determinant is \((\det A)^2\).
For singular \(A\), the left side is zero and the right side nonnegative.
(c) In the nonsingular case, equality means \(AA^\top\) diagonal, or
mutually orthogonal rows. In the singular case equality means at least
one row squared norm is zero, or at least one zero row. Thus equality
holds exactly when a row is zero or all rows are mutually orthogonal.
\end{solution}

\begin{problem}\label{ex:12:ineq-ma-27}
(a) For real symmetric matrices, show that
\(\det A=\prod_i a_{ii}\iff A\text{ is diagonal}\) remains true in
order two, but can fail in every order at least three.
(b) Prove that a real symmetric matrix is diagonal exactly when its
multiset of eigenvalues equals its multiset of diagonal entries.
\end{problem}
\begin{solution}
(a) For \(A=\begin{pmatrix}a&b\\b&d\end{pmatrix}\), the equality is
\(ad-b^2=ad\), equivalent to \(b=0\).
In order three, \(\begin{pmatrix}0&1&0\\1&0&0\\0&0&0\end{pmatrix}\)
is nondiagonal but both determinant and diagonal product vanish.
Appending an identity block gives the same counterexample in every larger order.
(b) If the multisets agree, then
\[
\sum_i a_{ii}^2=\sum_i\lambda_i^2=\tr(A^2)
=\sum_i a_{ii}^2+2\sum_{i<j}a_{ij}^2.
\]
Thus every off-diagonal entry vanishes. Conversely a diagonal matrix
has exactly its diagonal entries as its eigenvalues.
\end{solution}

\needspace{8\baselineskip}
\subsection*{Inverse and trace comparisons}

\begin{problem}\label{ex:12:ineq-ma-28}
(a) If \(A>0\), show \(A+A^{-1}-2I\geqslant0\).
(b) When is this matrix positive definite?
(c) For positive definite \(A,B\), prove
\(\tr((A^{-1}-B^{-1})(A-B))\leqslant0\).
\end{problem}
\begin{solution}
(a,b) In an eigenbasis of \(A\), the entries are
\(\lambda+\lambda^{-1}-2=(\lambda-1)^2/\lambda\geqslant0\).
They are all positive exactly when \(1\) is not an eigenvalue of \(A\).
(c) Put \(C=A^{-1/2}BA^{-1/2}>0\). Cyclicity of trace gives
\[
\tr((A^{-1}-B^{-1})(A-B))
=2n-\tr(A^{-1}B)-\tr(B^{-1}A)
=-\tr(C+C^{-1}-2I)\leqslant0
\]
by part (a).
\end{solution}

\begin{problem}\label{ex:12:ineq-ma-30}
If \(A\) and \(A+B\) are positive definite, prove
\[
\frac{\det(A+B)}{\det A}\leqslant\exp(\tr(A^{-1}B)),
\]
with equality exactly when \(B=0\).
\end{problem}
\begin{solution}
Set \(C=A^{-1/2}(A+B)A^{-1/2}>0\).
The log-determinant inequality gives
\[
\log\frac{\det(A+B)}{\det A}=\log\det C
\leqslant\tr C-n=\tr(A^{-1}B).
\]
Exponentiation preserves the inequality. Equality is equivalent to
\(C=I\), equivalent to \(B=0\).
\end{solution}

\begin{problem}\label{ex:12:ineq-ma-31}
Let \(A>0\), let \(B\) be symmetric, and assume \(C=A+B\) is nonsingular.
Prove
\(C^{-1}BC^{-1}\leqslant A^{-1}-C^{-1}\).
Show that the inequality is strict exactly when \(B\) is nonsingular.
\end{problem}
\begin{solution}
Multiplying the proposed difference by \(C\) on both sides gives
\[
C(A^{-1}-C^{-1}-C^{-1}BC^{-1})C
=CA^{-1}C-C-B=BA^{-1}B.
\]
Thus the difference is
\(C^{-1}BA^{-1}BC^{-1}=(A^{-1/2}BC^{-1})^\top(A^{-1/2}BC^{-1})\),
which is positive semidefinite. It is positive definite exactly when
\(A^{-1/2}BC^{-1}\) has trivial kernel, exactly when \(B\) is nonsingular.
No positive-definiteness assumption on \(C\) is required.
\end{solution}

\needspace{8\baselineskip}
\subsection*{Block and entrywise products}

\begin{problem}\label{ex:12:ineq-ma-32}
Define the Hadamard product by \((A\odot B)_{ij}=a_{ij}b_{ij}\).
(a) For \(\lambda=(\lambda_i)\), \(\Lambda=\operatorname{diag}(\lambda_i)\),
prove \(\lambda^\top(A\odot B)\lambda=\tr(A\Lambda B^\top\Lambda)\).
(b) If \(A,B>0\), prove \(A\odot B>0\).
\end{problem}
\begin{solution}
(a) The trace is \(\sum_{i,j}a_{ij}\lambda_j b_{ij}\lambda_i\),
exactly the left-hand quadratic expression.
(b) For nonzero \(\lambda\), \(D=\Lambda B\Lambda\) is positive
semidefinite and nonzero: if \(\lambda_i\ne0\), its diagonal entry
\(\lambda_i^2b_{ii}>0\). Consequently \(A^{1/2}DA^{1/2}\) is positive
semidefinite and nonzero, so its eigenvalues have positive sum.
By (a) and trace cyclicity,
\(\lambda^\top(A\odot B)\lambda=\tr(AD)
=\tr(A^{1/2}DA^{1/2})>0\).
The entrywise product is symmetric, completing positive definiteness.
\end{solution}

\begin{problem}\label{ex:12:ineq-ma-33}
Let \(A>0\), \(B\in\R^{n\times m}\), and symmetric \(X\in\R^{m\times m}\).
Prove
\[
\begin{pmatrix}A&B\\B^\top&X\end{pmatrix}\geqslant0
\iff X-B^\top A^{-1}B\geqslant0,
\]
and prove the corresponding equivalence with both inequalities strict.
\end{problem}
\begin{solution}
For any \(u,v\), complete the square:
\[
u^\top Au+2u^\top Bv+v^\top Xv
=(u+A^{-1}Bv)^\top A(u+A^{-1}Bv)
+v^\top(X-B^\top A^{-1}B)v.
\]
If the last matrix is positive semidefinite, this proves the forward
quadratic bound. Conversely choose \(u=-A^{-1}Bv\) to isolate that matrix.
For strict positivity, if \(v\ne0\) the second term is positive;
if \(v=0,u\ne0\) the first is positive. Conversely the same substitution
with nonzero \(v\) forces strict positivity of the last matrix.
\end{solution}

\begin{problem}\label{ex:12:ineq-ma-34}
Let \(Z=\begin{pmatrix}A&B\\B^\top&D\end{pmatrix}>0\).
(a) Prove \(\det Z\leqslant\det A\det D\), with equality exactly at \(B=0\).
(b) If \(B\) is square, prove \((\det B)^2<\det A\det D\).
\end{problem}
\begin{solution}
The diagonal blocks are positive definite, and the preceding criterion
makes \(S=D-B^\top A^{-1}B>0\).
Congruence by \(\begin{pmatrix}I&-A^{-1}B\\0&I\end{pmatrix}\), of
determinant one, gives \(\det Z=\det A\det S\).
Since \(0<S\leqslant D\), determinant monotonicity gives (a), with
equality exactly when \(S=D\), or \(B^\top A^{-1}B=0\).
The latter is a Gram matrix, and vanishes exactly for \(B=0\).
(b) Put \(C=B^\top A^{-1}B\), so \(0\leqslant C<D\).
If \(C\) is singular, \(\det C=0<\det D\). Otherwise congruence by
\(C^{-1/2}\) gives eigenvalues strictly above one, hence again
\(\det C<\det D\). Now \(\det C=(\det B)^2/\det A\), proving the claim.
\end{solution}

\begin{problem}\label{ex:12:ineq-ma-35}
For every real rectangular \(A\), prove
\(\begin{pmatrix}I&A\\A^\top&A^\top A\end{pmatrix}\geqslant0\).
\end{problem}
\begin{solution}
The matrix factors as
\(\begin{pmatrix}I\\A^\top\end{pmatrix}(I\mid A)\).
Its quadratic form at \((x,y)\) is therefore \(\norm{x+Ay}^2\geqslant0\).
The factorization also makes symmetry explicit.
\end{solution}

\begin{problem}\label{ex:12:ineq-ma-37}
For \(A,B\in\R^{m\times n}\), prove
\begin{enumerate}
\item \(Z=\begin{pmatrix}I_m+AA^\top&A+B\\(A+B)^\top&I_n+B^\top B\end{pmatrix}\geqslant0\);
\item \(I_n+B^\top B\geqslant(A+B)^\top(I_m+AA^\top)^{-1}(A+B)\);
\item if \(m=n\), then
\((\det(A+B))^2\leqslant\det(I+B^\top B)\det(I+AA^\top)\).
\end{enumerate}
\end{problem}
\begin{solution}
(a) Its quadratic form at \((x,y)\) is
\(\norm{x+By}^2+\norm{A^\top x+y}^2\), a sum of nonnegative terms.
(b) The first block is positive definite, since
\(x^\top(I+AA^\top)x=\norm{x}^2+\norm{A^\top x}^2>0\) for \(x\ne0\).
Apply the Schur-complement criterion to obtain the inequality.
(c) Both matrices in (b) are positive semidefinite, so determinant
monotonicity gives
\[
\frac{(\det(A+B))^2}{\det(I+AA^\top)}
\leqslant\det(I+B^\top B).
\]
The denominator is positive; multiplying proves the claim.
\end{solution}

\begin{problem}\label{ex:12:ineq-ma-38}
For positive definite \(A,B\) and symmetric \(X\) of the same order,
prove
\[
\begin{pmatrix}A+B&A\\A&A+X\end{pmatrix}\geqslant0
\iff X\geqslant-(A^{-1}+B^{-1})^{-1}.
\]
\end{problem}
\begin{solution}
The Schur-complement criterion makes the left condition equivalent to
\(X\geqslant-A+A(A+B)^{-1}A\).
Let \(H=A-A(A+B)^{-1}A=A(A+B)^{-1}B\).
All three factors in the last expression are nonsingular, and
\[
H^{-1}=B^{-1}(A+B)A^{-1}=B^{-1}+A^{-1}.
\]
Thus \(H=(A^{-1}+B^{-1})^{-1}\), and substitution proves the desired
condition. This also establishes symmetry and positive definiteness of \(H\).
\end{solution}

\needspace{8\baselineskip}
\subsection*{Extreme eigenvalues and variation}

\begin{problem}\label{ex:12:ineq-ma-39}
For real symmetric \(A\) with eigenvalues \(\lambda_1\geqslant\cdots\geqslant\lambda_n\),
prove the Rayleigh bounds
\(\lambda_n\leqslant x^\top Ax/(x^\top x)\leqslant\lambda_1\) for \(x\ne0\).
Show the bounds are attained and express them as a minimum and maximum.
Explain why these variational descriptions are useful.
\end{problem}
\begin{solution}
In an orthonormal eigenbasis, write \(x=\sum_i c_iq_i\). Its quotient is
\(\sum_i\lambda_i c_i^2/\sum_i c_i^2\), a weighted average of the
eigenvalues. This proves both bounds. Taking \(x=q_1\), respectively
\(q_n\), attains them, so they are the maximum and minimum over nonzero
\(x\), equivalently over unit vectors. The representation permits bounds
and comparisons of extreme eigenvalues through quadratic forms, without
solving a characteristic equation; the next exercise demonstrates this.
\end{solution}

\begin{problem}\label{ex:12:ineq-ma-40}
For real symmetric \(A,B\), use the Rayleigh representations to prove
\[
\lambda_1(A+B)\leqslant\lambda_1(A)+\lambda_1(B),\qquad
\lambda_n(A+B)\geqslant\lambda_n(A)+\lambda_n(B).
\]
If \(B\geqslant0\), show also
\(\lambda_1(A+B)\geqslant\lambda_1(A)\) and
\(\lambda_n(A+B)\geqslant\lambda_n(A)\).
\end{problem}
\begin{solution}
For each unit \(x\),
\(x^\top(A+B)x\leqslant\lambda_1(A)+\lambda_1(B)\).
Taking its maximum gives the first bound. Bounding each summand below
and taking the minimum gives the second.
When \(B\geqslant0\), every unit \(x\) satisfies
\(x^\top(A+B)x\geqslant x^\top Ax\). Taking maxima proves monotonicity
of the largest eigenvalue, and taking minima proves monotonicity of the
smallest eigenvalue.
\end{solution}


\needspace{8\baselineskip}
\subsection*{Concavity of determinant functions}
\begin{problem}\label{ex:12:ineq-ma13-66}
A real-valued function on positive definite matrices is convex if
\(f(tA+(1-t)B)\leqslant tf(A)+(1-t)f(B)\) for \(0<t<1\), and concave
if the inequality is reversed. Prove (a) \(\log\det X\) is concave;
(b) for order \(n\geqslant2\), \(\det X\) is neither convex nor concave;
(c) \((\det X)^{1/n}\) is concave for \(n\geqslant2\).
\end{problem}
\begin{solution}
(a) The determinant inequality with positive weights gives
\(\det(tA+(1-t)B)\geqslant(\det A)^t(\det B)^{1-t}\).
All determinants are positive, so taking logarithms gives precisely
concavity of \(\log\det\).
(b) In order two take \(A=I_2\), \(B=\operatorname{diag}(\beta,5)\),
\(\beta>0\), and \(t=1/2\). Then
\[
\tfrac12\det A+\tfrac12\det B-\det((A+B)/2)=\beta-1.
\]
This difference is negative at \(\beta=1/2\), contradicting convexity,
and positive at \(\beta=2\), contradicting concavity.
Appending an identity block gives the same counterexamples in every
larger order. For order one the determinant is linear, so the dimension
restriction is essential.
(c) Apply the determinant-root inequality to \(tA,(1-t)B\):
\[
(\det(tA+(1-t)B))^{1/n}
\geqslant t(\det A)^{1/n}+(1-t)(\det B)^{1/n},
\]
since \(\det(tA)=t^n\det A\). This is the required concavity.
\end{solution}

